% 第一部分：引言与相关工作（由 main_zh.tex 引入）。

\section{引言}\label{sec:intro}

强化学习需要奖励，而大多数任务的奖励都得有人来写。MuJoCo 运动控制基准自带工程化的奖励：前进速度加上存活奖励，再减去控制代价，各项权重靠人工调出。如果一个语言模型能根据任务描述来评判行为，这一步就可以省掉。这个想法已有多种形式：语言模型编写奖励代码~\cite{yu2023language,ma2024eureka}；在文本环境中直接充当奖励函数~\cite{kwon2023reward}；视觉--语言模型为机器人姿态的图像打分~\cite{rocamonde2024vlm}；把语言模型对事件描述的偏好蒸馏成内在奖励~\cite{klissarov2024motif}。

本文研究这个问题的一个窄版本，并配齐回答它所需的对照。设想一个 40 亿参数的评判模型读一段文字摘要，内容是机器人在 32 步窗口内的表现。用它的分数训练 PPO，能否替代标准基准的工程化奖励？代价又是多少？我们比较同一家族、同样大小的两个评判模型。一个是 \kev{}~\cite{kev4b}，一种决策模型：对 JSON 状态提出的问题，它一次前向传播就给出答案，带指针头和校准温度。另一个是自回归（AR）评判模型 Qwen3.5-4B，以同样的方式读出它在各答案选项上的下一词元分布。两个评判模型看到的文本完全相同。另设一个对照条件 \texttt{env\_window}：用真实奖励，但时间分辨率与评判模型相同。它把“按窗口给反馈”的影响与“评判模型本身”的影响分开。

\paragraph{贡献。}
\begin{itemize}
\item 在五个 MuJoCo 任务上，对 PPO 的七种奖励来源做受控比较：工程化奖励；它的窗口平均版本；\kev{} 的分数，带或不带“健康”门控；由 \kev{} 在 10\% 窗口上标注、在线蒸馏出的奖励模型；以及 AR 评判模型的期望读出和采样读出。每个任务训练 100 万步、5 个种子（采样读出的 AR 评判为 3 个），统计方法依照 \citet{agarwal2021deep}（\S\ref{sec:results}）。
\item 一项离线一致性研究：在同一批窗口上，把每种评判信号与环境奖励联系起来（\S\ref{sec:agreement}）。
\item 实测代价：每个 run 被评判的窗口数、评判耗时和总耗时（\S\ref{sec:cost}）。
\item 一项更正。本文的早期版本（v1--v3）报告的结果并非由这些实验产生：训练中的“评判模型”其实是一个手写的仿真器状态函数，排行榜是常数（见上方声明和 \file{CORRECTION.md}）。本版的每个数字都由仓库中的脚本从仓库中的日志计算得出。
\end{itemize}

\section{相关工作}\label{sec:related}

\paragraph{语言模型作为奖励设计者和奖励函数。} Eureka~\cite{ma2024eureka} 和 Language to Rewards~\cite{yu2023language} 用语言模型编写或参数化奖励代码。这样策略的奖励仍然便宜，但设计工作转移到了模型的编程能力和外层搜索循环上。像 \citet{kwon2023reward} 在谈判等文本游戏中那样直接把模型当作奖励，不需要代码，但训练期间就要调用模型。VLM-RM~\cite{rocamonde2024vlm} 用 CLIP 为人形机器人渲染画面中的姿态打分。Motif~\cite{klissarov2024motif} 把语言模型对 NetHack 事件描述的偏好蒸馏成奖励网络，我们的 \texttt{kev\_distill} 条件与之类似。我们的评判模型读的是遥测数据的文字摘要而非像素。这样把“评判”与“感知”分开，代价是摘要需要人来设计。

\paragraph{LLM 评判。} 用语言模型给输出打分~\cite{zheng2023judging,liu2023geval}已是评测中的常规做法。作为奖励模型，它们已知可被利用~\cite{gao2023scaling,skalse2022defining}。生成式验证器~\cite{zhang2024genverifier}读取评判模型回答“是/否”的概率，我们的 AR 评判也是如此。\kev{} 的不同之处在于：它的答案是专用输出头给出的分布，经过温度校准，并且多个问题在一次前向传播中算完。

\paragraph{深度强化学习的评估。} 每个任务只有少数种子时，平均回报的点估计并不可靠。我们依照 \citet{agarwal2021deep}：对全部 run 计算四分位间均值（IQM），给出分层 bootstrap 置信区间，以及改进概率。各任务的回报也一并报告。
